\chapter{Methodology}
\label{sec:methodology}

This chapter details the framework for our task-conditioned predictive model. Building on the foundations of EBMs and JEPAs, we design an autoregressive ``imagination'' model. Given a current state and a text description of a task, this model predicts a physically coehrerent sequence of future latent states. By operating entirely in a compressed latent space rather than generating raw pixels, the model avoids heavy computational costs. At the same time, it learns an information-rich representation of the visual environment, preserving the exact semantic features a robot needs to plan and execute tasks.

The remainder of this chapter is structured as follows. \cref{sec:latent_encoders} defines the specialized neural encoders used to project high-dimensional visual observations and task descriptions into their corresponding latent representations. \cref{sec:latent_predictive_model} introduces the autoregressive predictive model responsible for simulating future latent trajectories. \cref{sec:training_ebm} details the energy-based formulation of the training objective, including the integration of a multi-step rollout loss for temporal coherence and Sketched Isotropic Gaussian Regularization (SIGReg) to mathematically prevent latent collapse. Finally, \cref{sec:model_architectures} specifies the concrete neural network architectures, including Vision Transformers and Parameter-Efficient Fine-Tuned (PEFT) foundation models, employed to realize this system.

\section{Latent Encoders}
\label{sec:latent_encoders}

To facilitate predictive modeling and energy estimation within a structured manifold, sensory observations and task specifications must be mapped to a joint latent space. This transformation is performed by specialized neural architectures referred to as \textit{encoders}. Specifically, we define a visual encoder $g_\theta: \mathcal{I} \to \mathbb{R}^d$ and a linguistic encoder $g_\phi: \Sigma^\ast \to \mathbb{R}^d$. These models map visual observations $o_t \in \mathcal{I}$ at time $t$ and task-specific textual descriptions $\tau \in \Sigma^\ast$ to $d$-dimensional latent representations $\mathbf{e}_t$ and $\mathbf{z}$ respectively, as formulated in \cref{eq:visual_encoder} and \cref{eq:task_encoder}:
\begin{subequations}
	\begin{equation}
		\mathbf{e}_t  = g_\theta(o_t)
		\label{eq:visual_encoder}
	\end{equation}
	\begin{equation}
		\mathbf{z}    = g_\phi(\tau)
		\label{eq:task_encoder}
	\end{equation}
\end{subequations}

\nomenclature{$\mathcal{I}$}{The manifold of all visual observations/images}
\nomenclature{$\Sigma^\ast$}{The set of all language-based task descriptions}
\nomenclature{$o_t$}{Visual observation at time $t$}
\nomenclature{$\tau$}{Textual description of a specific task}
\nomenclature{\( g_\theta \)}{Neural encoder for visual observations}
\nomenclature{\( g_\phi \)}{Neural encoder for task descriptions}
\nomenclature{$\mathbf{e}_t$}{Latent representation of $o_t$ in $\mathbb{R}^d$}
\nomenclature{$\mathbf{z}$}{Latent task representation in $\mathbb{R}^d$}

\section{Latent Predictive Model}
\label{sec:latent_predictive_model}

The dynamics of the environment are modeled in the latent space through a predictor $h_\theta: (\mathbb{R}^d \times \mathbb{R}^d) \to \mathbb{R}^d$. Conditioned on the current latent state \( \mathbf{e}_t \) and the task representation \( \mathbf{z} \), the predictor estimates the subsequent latent state \( \hat{\mathbf{e}}_{t+1} \) corresponding to a successful execution of task \( \tau \):
\begin{equation}
	\hat{\mathbf{e}}_{t+1} = h_\theta(\mathbf{e}_t, \mathbf{z}) .
	\label{eq:predictor}
\end{equation}

By iteratively applying $h_\theta$, we can generate autoregressive trajectories in the latent space. Starting from an initial observation \( o_t \), a latent sequence of \( T \) future states can be synthesized using \cref{algo:rollout}, providing a mechanism for long-horizon temporal reasoning.

\begin{algorithm}
	\caption{\textsc{Rollout}$(\mathbf{e}_t, \mathbf{z}, h_\theta, T)$}
	\label{algo:rollout}

	\SetKwProg{Fn}{Function}{}{}
	\Fn{\textsc{Rollout}$(\mathbf{e}_t, \mathbf{z}, h_\theta, T)$}{
		$\hat{\mathbf{e}}_{t} \gets \mathbf{e}_t$\;

		\For{$k \gets 0$ \KwTo $T-1$}{
			$\hat{\mathbf{e}}_{t+k+1} \gets h_\theta(\hat{\mathbf{e}}_{t+k}, \mathbf{z})$\;
		}

		\Return $\hat{\mathbf{e}}_{t+1:t+T}$\;
	}
\end{algorithm}

\nomenclature{\( h_\theta \)}{Task-conditioned latent transition model}
\nomenclature{\( \hat{\mathbf{e}}_{t+1} \)}{Predicted future latent state}
\nomenclature{\( a_{i:j} \)}{Sequence of quantity \( a \) from timestep \( i \) to \( j \)}

\section{Training the Energy-Based Model}
\label{sec:training_ebm}

This section provides a detailed overview of the energy-based formulation and process of training the latent-predictive model. \cref{sec:method_ebm_formulation} formulates the latent predictive model in terms of an EBM. Further, \cref{sec:method_pred_obj}, \cref{sec:method_reg_collapse_prev} and \cref{sec:ms_roll_phy} provide a breakdown of the different terms in the loss function, with \cref{sec:combined_optimization} formalizes the training of the model as an optimization on the parameters of the energy-based model.

\subsection{Energy-Based Model Formulation}
\label{sec:method_ebm_formulation}

The visual encoder and the latent predictor collectively define an Energy-Based Model (EBM), denoted as $f_\theta = \{g_\theta, h_\theta\}$. The energy associated with a transition from observation \( o_t \) to a candidate state \( o'_{t+1} \), given a task \( \tau \), is quantified by the $L_2$ distance between the predicted and encoded latents:
\begin{equation}
	E(o_t, o'_{t+1}, \tau) = \| h_\theta\left(g_\theta(o_t), g_\phi(\tau)\right) - g_\theta(o'_{t+1}) \|^2_2 ,
	\label{eq:loss_energy}
\end{equation}
Lower energy levels signify transitions with higher statistical likelihood. This relationship is formalized via the Boltzmann distribution, which models the unnormalized conditional probability:
\begin{equation}
	p(o'_{t+1} \mid o_t, \tau) = \frac{\exp(-E(o_t, o'_{t+1}, \tau))}{\int_{\tilde{o} \in \mathcal{I}}{\exp(-E(o_t, \tilde{o}, \tau))\ d\tilde{o}}}\ .
	\label{eq:boltzmann}
\end{equation}
In practice, the normalizing partition function in the denominator is intractable; however, it can be treated as a constant during optimization to facilitate maximum likelihood estimation.

\subsection{Predictive Objective}
\label{sec:method_pred_obj}

During training, we utilize trajectories sampled from agent episodes. For a ground-truth transition $\{ o_t, o_{t+1} \}$, the model objective is to minimize the energy, ideally reaching zero, which corresponds to $p(o_{t+1} \mid o_t, \tau) = 1$. The predictive loss is thus defined as:
\begin{subequations}
	\begin{align}
		\mathcal{L}_\mathrm{pred} & = E(o_t, o_{t+1}, \tau)                                                                                       \\
		\mathcal{L}_\mathrm{pred} & = \| h_\theta\left(g_\theta(o_t), g_\phi(\tau)\right) - g_\theta(o_{t+1}) \|^2_2 . \label{eq:predictive_loss}
	\end{align}
	\label{eq:pred_loss}
\end{subequations}

\subsection{Regularization and Latent Collapse Mitigation}
\label{sec:method_reg_collapse_prev}

A naive minimization of \( \mathcal{L}_\mathrm{pred} \) from \cref{eq:predictive_loss} poses a significant risk of \textit{latent collapse}. In this degenerate state, the encoder $g_\theta$ maps all observations to a constant vector $\mathbf{k}$, while the predictor $h_\theta$ reduces to an identity mapping. This results in zero loss but eliminates the information-carrying capacity of the latent space, rendering the representations indistinguishable.

To prevent this, we employ Sketched Isotropic Gaussian regularization (SIGReg)~\cite{lejepa}, which provides information-theoretic guarantees against collapse. As illustrated in \cref{fig:sigreg}, the method samples a batch of latents \( \{\mathbf{e}_i\}_{i=1}^{N} \) and projects them onto $M$ random directions \( \{\omega_j\}_{j=1}^{M} \) sampled from the unit sphere \( \mathcal{S}^d \):
\begin{equation}
	s_{ij} = \omega_j^\top \mathbf{e}_i ,
	\label{eq:projection}
\end{equation}
The Epps-Pulley test statistic is then computed between the scalar projections $\{s_{ij}\}_{i=1}^{N}$ and a standard normal distribution \( \mathcal{N}(0, 1) \). The SIGReg loss is the average statistic across all directions:
\begin{equation}
	\mathcal{L}_\mathrm{SIGReg} = \frac{1}{M} \sum_{j=1}^{M}{\mathrm{EP}\left(\{s_{ij}\}_{i=1}^{N}, \mathcal{N}(0, 1)\right)} .
	\label{eq:loss_sigreg}
\end{equation}

\begin{figure}
	\begin{center}
		\begin{tikzpicture}

	% Axes
	\draw[->] (-2,0) -- (2,0) node[right] {$e_1$};
	\draw[->] (0,-2) -- (0,2) node[above] {$e_2$};

	% Latent points
	\foreach \x/\y in {0.6/1.0, -0.9/0.8, 1.1/-0.6, -1.1/-1.0, 0.4/-0.8} {
			\fill[blue] (\x,\y) circle (2pt);
		}

	% Projection direction
	\draw[->, red, thick] (0,0) -- (1.4,1.2) node[right] {$\omega_j$};

	% Projection line (axis along v)
	\draw[thick, gray] (-1.5,-1.3) -- (1.5,1.3);

	% Projection lines
	\foreach \x/\y in {0.6/1.0, -0.9/0.8, 1.1/-0.6, -1.1/-1.0, 0.4/-0.8} {
			\draw[dashed] (\x,\y) -- ({(\x+\y)/2},{(\x+\y)/2});
		}

	\node at (0,-2.3) {\small $s_{ij} = \omega_j^\top \mathbf{e}_i$};

\end{tikzpicture}

	\end{center}
	\caption{Visualization of Sketched Isotropic Gaussian Regularization. Latents are projected onto random unit vectors to ensure the distribution remains non-degenerate and follows an isotropic Gaussian profile.}\label{fig:sigreg}
\end{figure}

\nomenclature{SIGReg}{Sketched Isotropic Gaussian Regularization}
\nomenclature{\( \mathcal{S}^d \)}{The $d$-dimensional unit sphere}
\nomenclature{\( \mathcal{N}(\mu, \sigma) \)}{Gaussian distribution with mean $\mu$ and variance $\sigma^2$}

\subsection{Multi-Step Rollout and Physical Coherence}
\label{sec:ms_roll_phy}

While the single-step predictive loss encourages local accuracy, it is insufficient for ensuring long-term physical coherence in autoregressive generation. To address compounding errors and enforce temporal consistency, we introduce a $K$-step rollout loss. Given an initial observation \( o_t \) and the ground-truth sequence \( o_{t+1:t+K} \), the rollout loss is defined as:
\begin{equation}
	\mathcal{L}_\mathrm{RO} = \sum_{i=1}^{K}{\| \hat{\mathbf{e}}_{t+i} - g_\theta(o_{t+i}) \|^2_2} ,
	\label{eq:loss_rollout}
\end{equation}
where \( \hat{\mathbf{e}}_{t+1:t+K} \) is obtained via the \textsc{Rollout} procedure (\cref{algo:rollout}). This objective ensures that the predictor remains tethered to physically plausible trajectories over extended horizons, as shown in \cref{fig:rollout}.

\begin{figure}
	\begin{center}
		\begin{tikzpicture}[
	% Increase horizontal and vertical distances
	node distance=1.5cm and 2.5cm,
	font=\sffamily\small,
	% Professional Palette from dvipsnames
	enc_style/.style={rectangle, draw=CornflowerBlue, thick, fill=CornflowerBlue!5, minimum width=1.6cm, minimum height=0.8cm, rounded corners=2pt},
	pred_style/.style={rectangle, draw=Goldenrod, thick, fill=Goldenrod!5, minimum width=1.6cm, minimum height=0.8cm, rounded corners=2pt},
	loss_style/.style={rectangle, draw=Gray, thick, fill=White, minimum width=1.6cm, minimum height=0.7cm, rounded corners=2pt},
	input_style/.style={font=\sffamily\bfseries},
	arrow/.style={-{Stealth[scale=1.2]}, thick, draw=Black!80}
	]

	% --- 1. TASK ENCODING (Top Left) ---
	\node [input_style] (tau) at (0, 1.8) {$\tau$};
	\node [enc_style] (g_phi) at (1.8, 1.8) {$g_\phi$};
	\node [right=0.6cm of g_phi, font=\bfseries] (z) {$\mathbf{z}$};

	% --- 2. THE STARTING POINT (Left) ---
	\node [input_style] (ot) at (0, 0) {$o_t$};
	\node [enc_style] (enc_t) at (1.8, 0) {$g_\theta$};

	% --- 3. THE ROLLOUT CHAIN (Predictors - Spread Horizontally) ---
	% Increased spacing between steps to 3.5cm
	\node [pred_style] (h1) at (4.3, 0) {$h_\theta$};
	\node [pred_style] (h2) at (7.3, 0) {$h_\theta$};
	\node [right=0.6cm of h2, font=\large] (dots) {$\dots$};
	\node [pred_style] (hK) at (11.3, 0) {$h_\theta$};

	% --- 4. GROUND TRUTH ENCODERS (Bottom) ---
	\node [enc_style] (gt1) at (4.3, -2.5) {$g_\theta$};
	\node [enc_style] (gt2) at (7.3, -2.5) {$g_\theta$};
	\node [enc_style] (gtK) at (11.3, -2.5) {$g_\theta$};

	% Ground Truth Inputs
	\node [input_style, below=0.6cm of gt1] (ot1) {$o_{t+1}$};
	\node [input_style, below=0.6cm of gt2] (ot2) {$o_{t+2}$};
	\node [input_style, below=0.6cm of gtK] (otK) {$o_{t+K}$};

	% --- 5. PER-STEP SQUARED ERRORS (Centered between Chain and GT) ---
	\node [loss_style] (d1) at (4.3, -1.25) {$\|\cdot\|^2$};
	\node [loss_style] (d2) at (7.3, -1.25) {$\|\cdot\|^2$};
	\node [loss_style] (dK) at (11.3, -1.25) {$\|\cdot\|^2$};

	% --- 6. FINAL AGGREGATION (Far Right) ---
	\node [rectangle, draw=BrickRed, fill=BrickRed!5, thick, minimum width=1.2cm, minimum height=1cm, rounded corners=3pt]
	(final_loss) at (13.5, -1.25) {$\displaystyle \mathcal{L}_{\mathrm{RO}}$};

	% --- CONNECTIONS ---

	% Task Flow (Clean horizontal line with vertical drops)
	\draw [arrow] (tau) -- (g_phi);
	\draw [arrow] (g_phi) -- (z);
	\coordinate (z_line) at (z.center);
	\draw [arrow, Black!40] (z.east) -| (h1.north);
	\draw [arrow, Black!40] (z.east) -| (h2.north);
	\draw [arrow, Black!40] (z.east) -| (hK.north);

	% Autoregressive Flow (With clear labels)
	\draw [arrow] (ot) -- (enc_t);
	\draw [arrow] (enc_t) -- node[above=0.1cm] {$\mathbf{e}_t$} (h1);
	\draw [arrow] (h1) -- node[above=0.1cm] {$\hat{\mathbf{e}}_{t+1}$} (h2);
	\draw [arrow] (h2) -- (dots);
	\draw [arrow] (dots) -- (hK);
	\draw [arrow] (hK) -- node[pos=1.5] {$\hat{\mathbf{e}}_{t+K}$} ++(1.7,0);

	% Ground Truth Flow
	\draw [arrow] (ot1) -- (gt1);
	\draw [arrow] (ot2) -- (gt2);
	\draw [arrow] (otK) -- (gtK);

	% Comparison Connections (Vertical comparisons)
	\draw [arrow, Gray] (gt1) -- (d1);
	\draw [arrow, Gray] (h1) -- (d1);
	\draw [arrow, Gray] (gt2) -- (d2);
	\draw [arrow, Gray] (h2) -- (d2);
	\draw [arrow, Gray] (gtK) -- (dK);
	\draw [arrow, Gray] (hK) -- (dK);

	% Accumulation to Final Loss (Dashed line spanning the rollout)
	\draw [arrow, dashed, BrickRed!30, line width=0.8pt] (d1) -- (d2);
	\draw [arrow, dashed, BrickRed!30, line width=0.8pt] (d2) -- (dK);
	\draw [arrow, dashed, BrickRed!30, line width=0.8pt] (dK) -- (final_loss.west);

\end{tikzpicture}

	\end{center}
	\caption{\textbf{Rollout Loss mechanism.}\; The loss is computed by comparing the sequence of latents generated by the predictor's autoregressive rollout against the latents encoded from the ground-truth visual observations.}\label{fig:rollout}
\end{figure}

\subsection{Combined Optimization Objective}
\label{sec:combined_optimization}

\begin{figure}
	\begin{center}
		\begin{tikzpicture}[
		node distance=1.2cm and 2.0cm,
		font=\sffamily\small,
		% Styles
		block/.style={rectangle, draw, thick, minimum width=1.8cm, minimum height=1cm, align=center, fill=white, rounded corners=2pt},
		loss/.style={block, fill=gray!5, minimum width=2.2cm},
		pred_style/.style={block, fill=Goldenrod!10, draw=Goldenrod!60},
		enc_style/.style={block, fill=CornflowerBlue!10, draw=CornflowerBlue!60},
		input/.style={font=\sffamily\bfseries},
		arrow/.style={-Stealth, thick, line width=1pt},
		% Style for the EBM container
		ebm_box/.style={draw=PineGreen, dashed, line width=1.5pt, fill=PineGreen!5, fill opacity=0.3, rounded corners=5pt}
	]

	% --- Top Branch (Regularization) ---
	\node [block, label={[font=\bfseries]above:Buffer}] (buffer) {$\mathcal{B}$};
	\node [enc_style, right=of buffer] (reg_enc) {$g_\theta$};
	\node [loss, right=of reg_enc] (sig_loss) {$\mathcal{L}_\mathrm{SIGReg}$};
	\node [above=-0.1cm of sig_loss, font=\scriptsize\itshape] {Regularization};

	% --- Middle Branch (Prediction) ---
	% Position enc_t below the top branch
	\node [enc_style, below=1.5cm of reg_enc] (enc_t) {$g_\theta$};
	\node [pred_style, right=of enc_t] (pred) {${h_\theta}$};
	\node [loss, right=of pred] (l2_loss) {$\mathcal{L}_{\mathrm{pred}}$};

	% --- Bottom Branch (Target) ---
	\node [enc_style, below=1.5cm of pred] (enc_tp1) {$g_\theta$};

	% --- Inputs (Vertically Aligned) ---
	% ot aligned with enc_t
	\node [input, color=Magenta, left=2.6cm of enc_t] (ot) {$o_t$};
	% ot+1 forced to align vertically with ot
	\node [input, color=Blue] at (ot |- enc_tp1) (otp1) {$o_{t+1}$};
	\node [input, above=0.6cm of pred, color=RedOrange] (z) {$z$ (Task)};

	% --- EBM f_theta Container ---
	% This box wraps the components that make up the EBM
	\begin{scope}[on background layer]
		\node [ebm_box, fit=(enc_t) (enc_tp1) (l2_loss) (pred), label={[PineGreen, font=\bfseries\large]below:EBM $f_\theta$}] (ebm_container) {};
	\end{scope}

	% --- Connections ---

	% Top path
	\draw [arrow] (buffer) -- node[above, font=\scriptsize] {$b_0 \sim \mathcal{B}$} (reg_enc);
	\draw [arrow] (reg_enc) -- (sig_loss);

	% Middle path
	\draw [arrow, color=Magenta] (ot) -- (enc_t);
	\draw [arrow, dashed, color=Magenta] (ot) -- (buffer);

	\draw [arrow] (enc_t) -- node[above, color=Magenta!80!black, font=\scriptsize] {$e_t$} (pred);
	\draw [arrow, color=RedOrange] (z) -- (pred);
	\draw [arrow] (pred) -- node[above, color=Blue!80!black, font=\scriptsize] {$\hat{e}_{t+1}$} (l2_loss);

	% Bottom path
	\draw [arrow, color=Blue] (otp1) -- (enc_tp1);
	\draw [arrow, color=Blue] (enc_tp1) -| node[pos=0.25, above, font=\scriptsize] {$e_{t+1}$} (l2_loss);

	% Energy Output Arrow
	\draw [arrow, draw=PineGreen, line width=1.5pt] (l2_loss.north) -- ++(0,1.5)
	node[above, PineGreen, font=\bfseries] {$E(\textcolor{Magenta}{o_t}, \textcolor{Blue}{o_{t+1}}, \textcolor{RedOrange}{z})$};

\end{tikzpicture}

	\end{center}
	\caption{\textbf{Schematic of EBM training pipeline.} The encoder $g_\theta$ and predictor $h_\theta$ learn to predict the latent of the observation at the next timestep $\hat{\mathbf{e}}_{t+1}$, given the current observation $o_t$ and task $\tau$. The SIGReg loss, calculated using a separate buffering and sampling mechanism, prevents collapse of the latent space of the encoder.}\label{fig:architecture}
\end{figure}

The comprehensive training objective combines predictive accuracy, temporal coherence, and latent space entropy. To ensure that \( \mathcal{L}_\mathrm{SIGReg} \) captures the global distribution of the latent space rather than just the local trajectory, we maintain a replay buffer \( \mathcal{B} \). Batches of observations \( b_0 \sim \mathcal{B} \) are sampled to compute the regularization term, as depicted in \cref{fig:architecture}. The final weighted loss function is:
\begin{equation}
	\mathcal{L} = \mathcal{L}_\mathrm{pred}
	+ \lambda_\mathrm{RO} \mathcal{L}_\mathrm{RO}
	+ \lambda_\mathrm{SIGReg} \mathcal{L}_\mathrm{SIGReg} ,
	\label{eq:loss_combined}
\end{equation}
where \( \lambda_\mathrm{RO} \) and \( \lambda_\mathrm{SIGReg} \) are hyperparameters regulating the contribution of the rollout and regularization terms, respectively. The \textit{training} of the model is defined as an optimization, formalized as
\begin{equation}
	(\theta^\ast, \gamma^\ast) = \underset{(\theta, \gamma)}{\arg \min\ }{\mathcal{L}(\theta, \gamma)} ,
	\label{eq:training_optimization}
\end{equation}
where \( \mathcal{L}(\theta, \gamma) \) specifies that the combined loss \( \mathcal{L} \) is a function of the parameters of the encoder, predictor and task encoder. Here, \( (\theta^ast, \gamma^\ast) \) is the combination of parameters that minimizes the combined loss function defined in \cref{eq:loss_combined}.

\nomenclature{\( \lambda_\mathrm{RO} \)}{Weighting hyperparameter for the rollout loss}
\nomenclature{\( \lambda_\mathrm{SIGReg} \)}{Weighting hyperparameter for the SIGReg loss}

\section{Model Architectures}
\label{sec:model_architectures}

\paragraph{Vision Encoder} We employ a Vision Transformer (\texttt{ViT})~\cite{vit} as the primary visual backbone. Our experiments demonstrate that a \texttt{ViT-Tiny} configuration (5M parameters) provides a sufficient balance between representational capacity and computational overhead.

\paragraph{Text Encoder} Task descriptions are encoded using CLIP~\cite{clip}, a multimodal foundation model. We specifically utilize the \texttt{clip-vit-base} variant. To investigate the benefits of task-specific semantic alignment, we include ablation studies involving parameter-efficient fine-tuning (PEFT)~\cite{hu2022lora} of the text encoder.

\paragraph{Predictive Backbone} We evaluate two distinct architectural paradigms for the predictor $h_\theta$: a Transformer-decoder~\cite{vaswani2017attention}, which leverages self-attention for sequence modeling, and a deep Multi-Layer Perceptron (MLP) utilizing GELU non-linearities.
